% ============================================================================
% CAPÍTULO 25 — Carreira\index{carreira}: do aprendizado ao mercado
% Status: texto completo (rodada editorial 2)
% ============================================================================
\chapter{Carreira: do Aprendizado ao Mercado}
\label{cap:carreira}

\begin{caixaconceito}
Você construiu o sistema inteiro — frontend, backend, dados, infraestrutura,
segurança, IA. Agora o desafio final: transformar esse conhecimento em
\emph{oportunidades}. Este capítulo é um plano de carreira prático, do
portfólio\index{portfólio} à entrevista técnica.
\end{caixaconceito}

Este é o último capítulo técnico — e talvez o mais importante. Você tem agora
o que pouquíssimas pessoas têm: a capacidade de construir um produto full
stack \emph{completo}, com testes, infraestrutura, segurança e IA, e de
explicar cada decisão pelo caminho. O mercado paga por essa capacidade — mas
não por osmose. Ele paga para quem \emph{mostra}. Este capítulo é o manual de
como mostrar.

Ao final deste capítulo, você será capaz de:
\begin{objetivos}
  \item Construir um portfólio e um currículo que contam uma história técnica;
  \item Posicionar-se no mercado (LinkedIn\index{LinkedIn}, comunidades, marca pessoal);
  \item Se preparar para entrevistas\index{entrevistas} técnicas (estrutura e prática);
  \item Entender os níveis de carreira e como negociar salário;
  \item Planejar o aprendizado contínuo em um mercado que muda rápido;
  \item Entregar o projeto do capítulo: seu plano de carreira individual.
\end{objetivos}

\section{O que o mercado procura em 2026}

\begin{itemize}[leftmargin=*, itemsep=0.4em]
  \item \textbf{Fundamentos sólidos} (HTTP, SQL, algoritmos, estrutura de dados);
  \item \textbf{Entrega de ponta a ponta}: não basta saber — precisa \emph{mostrar} (portfólio, projetos reais);
  \item \textbf{Qualidade}: testes, revisão, documentação — o que este livro treinou;
  \item \textbf{IA como ferramenta}: produtividade com assistentes, e capacidade de avaliar o que geram;
  \item \textbf{Comunicação}: explicar decisões, trade-offs e erros com clareza.
\end{itemize}

\begin{caixaconceito}
Note o que \emph{não} está na lista: ``saber 15 frameworks'' não está, ``ter
decorado a documentação'' não está. O mercado de 2026 compra três coisas:
\emph{fundamentos} (que nunca mudam de valor), \emph{entrega comprovada}
(projetos que funcionam e são explicáveis) e \emph{julgamento} (a capacidade
de decidir entre alternativas — exatamente o que este livro treinou em cada
capítulo).
\end{caixaconceito}

\section{O portfólio que funciona}

\begin{itemize}[leftmargin=*, itemsep=0.4em]
  \item \textbf{Qualidade $>$ quantidade}: 2--3 projetos completos valem mais que 10 incompletos;
  \item \textbf{Projetos que contam história}: o SkillHub (full stack), o blog (Next.js/ISR), o jogo (algoritmos) — cada um prova uma habilidade;
  \item \textbf{README de nível profissional}: problema, solução, arquitetura (diagrama), como rodar, testes, deploy;
  \item \textbf{Código limpo e testado}: recrutadores e seniors \emph{leem} o código;
  \item \textbf{Presença pública}: deploy funcional + link no GitHub + descrição no LinkedIn.
\end{itemize}

\begin{caixadica}
O teste definitivo do portfólio: \emph{se um recrutador abrir o repositório às
23h, ele entende o projeto em 5 minutos?} O README responde (problema,
arquitetura, como rodar), os testes provam (roda e passa), o deploy
demonstra (está no ar). Se alguma dessas três faltar, o projeto não está
pronto para o mercado — só para você.
\end{caixadica}

\section{Currículo e LinkedIn}

\begin{itemize}[leftmargin=*, itemsep=0.4em]
  \item \textbf{Formato} (recomendado em 2026): uma página, ordem inversa, foco em impacto (\emph{``reduzi o tempo de carregamento em 40\%''}, não \emph{``usei React''});
  \item \textbf{Métricas}: toda conquista com número mensurável;
  \item \textbf{LinkedIn}: título que descreve o que você \emph{entrega}, resumo com história, projetos fixados, recomendações de quem trabalhou com você;
  \item \textbf{Comunidades}: responder perguntas (Stack Overflow, fóruns) gera reputação técnica real.
\end{itemize}

\begin{caixaconceito}
A diferença entre ``usei React e Next.js'' e ``entreguei um marketplace com
pagamento, testes e deploy em produção'' não é de vocabulário — é de
\emph{evidência}. O currículo moderno não lista ferramentas; ele lista
resultados, e as ferramentas aparecem como os meios. Toda linha do seu
currículo deveria passar no teste: ``e daí?''.
\end{caixaconceito}

\section{Entrevistas técnicas: a estrutura}

\begin{enumerate}[leftmargin=*, itemsep=0.4em]
  \item \textbf{Screening} (RH): sua história, motivação, fit — prepare 30 segundos de apresentação;
  \item \textbf{Técnica 1 — algoritmos}: resolva problemas com pensamento em voz alta (complexidade, trade-offs);
  \item \textbf{Técnica 2 — sistema/domínio}: desenhe um sistema (``projete um encurtador de URLs'' — você já construiu um!);
  \item \textbf{Behavioral}: conte com método STAR (Situação, Tarefa, Ação, Resultado);
  \item \textbf{Cultural}: pergunte você também — time, processo, desafios.
\end{enumerate}

\begin{caixadica}
Seus projetos deste livro são o melhor material de entrevista: cada um já é um
\emph{case} com arquitetura, testes e decisões documentadas. ``Conte como você
decidiu entre Prisma e Drizzle'' é exatamente o tipo de pergunta que separa
quem fez de quem decorou.
\end{caixadica}

\section{Níveis, salários e negociação}

\begin{itemize}[leftmargin=*, itemsep=0.4em]
  \item \textbf{Júnior $\rightarrow$ Pleno $\rightarrow$ Sênior}: autonomia, escopo e impacto crescem — o salário segue o impacto, não os anos;
  \item \textbf{Pesquise o mercado} (Glassdoor, Levels.fyi, pesquisas salariais BR) antes de negociar;
  \item \textbf{Negocie o pacote todo}: salário, benefícios, equity, modelo híbrido/remoto, crescimento;
  \item \textbf{Nunca aceite a primeira oferta sem pensar}: peça tempo para avaliar;
  \item \textbf{Cresça por decisão, não por acaso}: peça feedback, assuma escopo maior, documente resultados.
\end{itemize}

\begin{caixaconceito}
A hierarquia de carreira não mede \emph{anos}, mede \emph{autonomia e
impacto}: o júnior executa tarefas bem definidas, o pleno resolve problemas
sem supervisão, o sênior define a direção técnica e desbloqueia o time. Se
você quer crescer, cresça o escopo do que assume e documente o resultado —
o título segue o impacto, não o contrário.
\end{caixaconceito}

\section{Aprendizado contínuo: o mercado muda a cada ano}

\begin{itemize}[leftmargin=*, itemsep=0.4em]
  \item \textbf{Sistema de aprendizado}: 70\% prática no trabalho/projetos, 20\% interação (comunidades, code review), 10\% estudo formal (este livro!);
  \item \textbf{Acompanhe fontes oficiais}: blogs do React/Next.js/Node, Octoverse, State of JS, Technology Radar — as mesmas que embasam este livro;
  \item \textbf{Profundidade primeiro}: dominar um ecossistema (TS/React/Next) vale mais que ``saber um pouco de tudo'';
  \item \textbf{IA como colega}: delegue o mecânico, mantenha o crítico (você treinou isso no capítulo~\ref{cap:ia});
  \item \textbf{Ensine}: blogar ou apresentar o que aprendeu consolida e constrói reputação.
\end{itemize}

\begin{caixadica}
A habilidade mais subestimada do mercado é \textbf{ensinar}. Explicar o que
você aprendeu — num artigo, numa palestra, num code review — força você a
organizar o conhecimento, revela as lacunas e constrói a reputação que
currículo nenhum compra. Quem ensina também aprende duas vezes.
\end{caixadica}

% ============================================================================
\section{Projeto do capítulo: seu plano de carreira individual}
\label{sec:projeto25}

\begin{caixaprojeto}
\textbf{Objetivo:} transformar tudo em ação: portfólio finalizado, currículo
atualizado e um plano de 12 meses documentado.

\textbf{Critérios de aceite:}
\begin{itemize}[leftmargin=*, itemsep=0.2em]
  \item Portfólio com 3 projetos completos (READMEs profissionais e deploys no ar);
  \item Currículo de 1 página com métricas de impacto;
  \item LinkedIn atualizado com projetos fixados e resumo com história;
  \item Plano de 12 meses: metas trimestrais, stack a aprofundar, ritmo de estudo (X horas/semana);
  \item Lista de 10 vagas-alvo com requisitos mapeados contra suas habilidades;
  \item Plano de estudo de algoritmos (2 problemas/semana, estruturas de dados);
  \item Respostas STAR escritas para as 5 perguntas comportamentais mais comuns.
\end{itemize}
\end{caixaprojeto}

\subsection{Passo a passo}

\begin{passos}
  \item \textbf{Inventário}: liste seus projetos e traduza cada um em uma
  habilidade comprovada (com métricas);
  \item \textbf{Portfólio}: escolha os 3 melhores projetos, escreva READMEs
  profissionais e garanta que todos tenham deploy no ar;
  \item \textbf{Currículo}: reescreva em uma página com foco em impacto;
  \item \textbf{LinkedIn}: atualize título, resumo e fixe os projetos;
  \item \textbf{Pesquisa de vagas}: liste 10 vagas-alvo e mapeie os requisitos
  contra o que você tem (e o que falta);
  \item \textbf{Plano de 12 meses}: defina metas trimestrais, horário de
  estudo e o plano de algoritmos;
  \item \textbf{Prática de entrevista}: escreva as respostas STAR e faça a
  primeira entrevista simulada cronometrada.
\end{passos}

\section{Exercícios de fixação}

\begin{exercicios}
  \item Escreva sua apresentação de 30 segundos (quem você é, o que entrega, o que busca).
  \item Liste as 5 habilidades técnicas mais pedidas nas vagas que você encontrou e onde você está em cada uma (1--5).
  \item Transforme uma conquista sua em uma frase com métrica.
  \item Explique o método STAR com um exemplo seu.
  \item Defina seu horário de estudo semanal realista (e o que será cortado para caber).
\end{exercicios}

\section{Desafios}

\begin{desafios}
  \item \textbf{Entrevista simulada}: resolva ``projete um encurtador de URLs'' (ou ``design do Instagram'') em voz alta, cronometrado, e grave para assistir depois.
  \item \textbf{Case de sistema}: desenhe o SkillHub em 15 minutos em um quadro branco (papel), cobrindo arquitetura, banco, cache, filas e segurança.
  \item \textbf{Marca pessoal}: escreva e publique um artigo técnico sobre algo que você construiu neste livro (pode ser no seu blog MDX!).
\end{desafios}

\section{Exercícios de nível sênior}

\begin{seniores}
  \item \textbf{Simulação de code review}: pegue um projeto público no GitHub, faça um code review completo (comentários por linha, priorizados) e abra o PR real.
  \item \textbf{Técnica de negociação}: simule uma negociação salarial com um colega, com contra-ofertas e ``pacote inteiro'', e documente o que aprendeu.
  \item \textbf{Mentoria}: encontre um iniciante e explique o fluxo de uma requisição HTTP até o banco e de volta — ensinar é o melhor teste de domínio.
  \item \textbf{Tendências}: leia o Octoverse e o Technology Radar mais recentes, escolha 3 tendências e monte seu plano de estudo para os próximos 2 anos.
\end{seniores}

\section{Armadilhas comuns}

\begin{itemize}[leftmargin=*, itemsep=0.4em]
  \item Estudar infinitamente sem aplicar (``vou começar quando estiver pronto'');
  \item Portfólio de tutoriais copiados (recrutadores reconhecem);
  \item Currículo sem métricas (``responsável por...'' vazio);
  \item Recusar feedback (a entrevista também serve para você aprender);
  \item Comparação com os outros em vez de comparação com o seu eu de ontem;
  \item Aplicar para tudo sem direção (a lista de vagas-alvo existe para isso).
\end{itemize}

\section{Resumo do capítulo}

\begin{itemize}[leftmargin=*, itemsep=0.3em]
  \item Portfólio com 3 projetos completos conta mais que 10 incompletos;
  \item Currículo e LinkedIn falam em impacto, com métricas;
  \item Entrevistas: screening, algoritmos, design de sistema, behavioral (STAR);
  \item Níveis seguem impacto; negocie o pacote todo;
  \item Aprendizado contínuo: prática + comunidade + fontes oficiais;
  \item Projeto entregue: seu plano de carreira de 12 meses.
\end{itemize}

\begin{caixaconceito}
\textbf{Epílogo.} No capítulo~1, você viu os dados de mercado que justificaram
esta jornada. No capítulo~25, você tem o sistema, o portfólio e o plano. O
que separa o leitor que chega até aqui do profissional que o mercado contrata
não é mais conhecimento — é \emph{constância}: aplicar o plano, medir o
progresso e repetir. O livro termina aqui; a sua jornada começa.
\end{caixaconceito}

\section{Referências oficiais para aprofundar}

\begin{itemize}[leftmargin=*, itemsep=0.3em]
  \item Levels.fyi (salários e níveis): \url{https://www.levels.fyi}
  \item Peneira: sistema de entrevistas (aberto): \url{https://peneira.dev}
  \item roadmap.sh (mapas de carreira): \url{https://roadmap.sh}
  \item GitHub Octoverse: \url{https://octoverse.github.com}
\end{itemize}
